\documentclass[11pt,a4paper]{article}
\usepackage[margin=1in]{geometry}
\usepackage{amsmath,booktabs,hyperref,float}

\title{Iteration 026: Gaussian Noise Augmentation}
\author{Autoresearch Loop}
\date{April 8, 2026}

\begin{document}
\maketitle

\begin{abstract}
SNR-preserving Gaussian noise augmentation with InstanceNorm achieves $r=0.375$,
not improving over the baseline ($r=0.378$). The InstanceNorm component degrades
performance, consistent with iterations 020--025.
\end{abstract}

\section{Hypothesis}
Measurement noise varies across subjects and electrodes. Adding channel-wise
Gaussian noise ($\sigma = 0.1 \times \text{channel std}$) during training
makes the model robust to these differences.

\section{Method}
During training, add Gaussian noise scaled to 10\% of each channel's
per-window standard deviation. Architecture: InstanceNorm + FIR (7 taps).
Combined loss + corr validation + 15\% channel dropout.

\section{Results}
\begin{table}[H]
\centering
\begin{tabular}{lrrr}
\toprule
Model & Mean $r$ & Std $r$ & SNR (dB) \\
\midrule
Combined loss (iter017) & 0.378 & 0.078 & 0.61 \\
Noise + InstanceNorm & 0.375 & 0.077 & 0.62 \\
\bottomrule
\end{tabular}
\end{table}

\section{Analysis}
Same InstanceNorm degradation pattern. The noise augmentation idea may have
merit without InstanceNorm, but was not separately tested in this round due
to the confounding InstanceNorm effect.

\end{document}
